% Candidate insertion after the unmerged higher-dimensional matching section.
% Frozen v0.3 and its existing revision section are not changed by this file.
% Requires the theorem environments and amsmath/amssymb used by the baseline.
\section{Interaction cores beyond matching}
\label{sec:interaction-cores}

Retain $K=\Delta(Q\times\mathbf 2)$, $L=\Delta Q^\psi$, the missing
sets $M_0=P\times\{0\}$ and $M_1=Z\times\{1\}$, and
$U=K_P\cup K_Z$.  Cone homology in this section is reduced homology
of the topological homotopy cofiber, equivalently homology of the algebraic
mapping cone.  The two-colour shift is
\[
 \widetilde H_n(\operatorname{Cone}(\Delta\pi);R)
 \cong H_{n+1}(K,U;R).
\]
For a mixed missing core $S$, put
$\Lambda_S=\operatorname{lk}_K(S)[V(L)]$.
All link chain complexes below are augmented, including degree $-1$.

\begin{proposition}[Interaction spectral sequence]
\label{prop:interaction-ss}
There is a finite convergent spectral sequence
\[
 E^1_{p,q}=\bigoplus_{|S|=p}\widetilde H_q(\Lambda_S;R)
 \quad\Longrightarrow\quad H_{p+q}(K,U;R),
\]
with $d^r$ of bidegree $(-r,r-1)$.  The first differential is the
signed sum of the inclusion maps obtained by deleting one missing vertex,
omitting nonmixed faces.  If all mixed cores have size at most $M\ge2$,
then $E^{M-1}=E^\infty$.
\end{proposition}
\begin{proof}
Orient a relative simplex by listing its $p$ missing vertices first and
its available vertices second, with the corresponding shuffle sign.
Its boundary is
\[
 d(S\wedge T)=\sum_i(-1)^i(S\setminus\{s_i\})\wedge T
                 +(-1)^pS\wedge\partial T.
\]
Terms with nonmixed core vanish in the relative complex.  Filtering by
$p$ gives the displayed page and differential.  The filtration is finite.
Only columns $2,\ldots,M$ can occur, so every $d^r$ with $r\ge M-1$
has zero source or target.
\end{proof}

\begin{theorem}[Exact two-column interaction theorem]
\label{thm:two-column-interaction}
Suppose no mixed core has four or more vertices.  Define
\[
 A_\bullet=\bigoplus_{|S|=2}\widetilde C_\bullet(\Lambda_S;R),
 \qquad
 B_\bullet=\bigoplus_{|S|=3}\widetilde C_\bullet(\Lambda_S;R),
\]
and let $h:B_\bullet\to A_\bullet$ be the signed sum of core-deletion
inclusions.  Then
\[
 C_\bullet(K,U;R)\cong\operatorname{Cone}(h)[2],\qquad
 \widetilde H_n(\operatorname{Cone}(\Delta\pi);R)
 \cong H_{n-1}(\operatorname{Cone}(h)).
\]
The interaction spectral sequence collapses at $E^2$, and there is a
natural exact sequence
\[
 0\longrightarrow\operatorname{coker}H_{n-1}(h)
 \longrightarrow\widetilde H_n(\operatorname{Cone}(\Delta\pi);R)
 \longrightarrow\ker H_{n-2}(h)\longrightarrow0.
\]
In particular the cone is $R$-acyclic if and only if $h$ is a
quasi-isomorphism.
\end{theorem}
\begin{proof}
The degree-$m$ relative group is $A_{m-2}\oplus B_{m-3}$ and its
boundary is $(a,b)\mapsto(\partial a+hb,-\partial b)$.
This is the asserted shifted chain cone.  Apply the two-colour shift,
the mapping-cone long exact sequence, and
Proposition~\ref{prop:interaction-ss} with $M=3$.
\end{proof}

The exact sequence splits noncanonically over a field.  Over $\mathbb Z$
it splits when the right-hand kernel is free abelian; no splitting is
asserted in general.  Thus collapse alone does not replace the integral
answer by a direct sum of link-homology kernels and cokernels.  Signed
incidence coefficients on link chains likewise do not imply total
unimodularity of the induced maps on homology.

\begin{theorem}[Degree-two classification and its sharp page bound]
\label{thm:degree-two-interaction}
Suppose the exclusive graph has maximum degree at most two.
Every mixed core has size at most four, so $E^3=E^\infty$.
If no $C_4$ component has both its same-colour pairs comparable in $Q$,
Theorem~\ref{thm:two-column-interaction} applies and $E^2=E^\infty$.
In every cycle component of length at least six the cone contribution is
the coarse-edge-link direct sum of the matching theorem.
\end{theorem}
\begin{proof}
Transitivity gives $N(p')\subseteq N(p)$ when $p<p'$ in $P$, and
$N(z)\subseteq N(z')$ when $z<z'$ in $Z$.
A mixed core is a biclique with each colour class a chain.  Degree two
bounds both class sizes by two.  A four-core requires a whole $K_{2,2}$
component with both colour pairs comparable.  This proves the bounds.
On a longer cycle, two same-colour degree-two vertices have distinct
neighbourhoods of the same cardinality and cannot be comparable if they
share a neighbour.  Hence there are no triple cores.  No missing vertex
can then lie in a product-edge link, so the original matching proof
applies componentwise without change.
\end{proof}

More explicitly, a path with at least three edges supports triples only
on its first two or last two edges: a comparable common-neighbour pair
must include a leaf.  For a path of at least four edges these end triples
have disjoint edge supports, giving independent star blocks and interior
edge blocks.  A three-edge path has up to two triples sharing its middle
edge and is treated by one instance of $\operatorname{Cone}(h)$.
Every surviving relative face remains in the same exclusive graph
component, so the component decomposition is at chain level.  Links
must nevertheless be calculated in the whole coarse poset.

For a two-edge star with ordered leaves, $h$ is the signed pair of
inclusions from its triple link to its two edge links.  On the frozen
three-chain profile $(1,0,0)$, all three available links have reduced
$H_0=\mathbb Z$ and $h_0$ is the primitive signed diagonal
$\mathbb Z\to\mathbb Z^2$.  This gives cone $H_1=\mathbb Z$,
although both coarse edge links are contractible.

\subsection{Higher differentials are necessary}
For a $C_4$ whose two pairs are comparable, write the core order
$a<b<x<y$.  Let $A,B,V$ be the sums of edge, triple, and four-core link
chains respectively.  Then
\[
 C_m(K,U)=A_{m-2}\oplus B_{m-3}\oplus V_{m-4},\qquad
 d(a,b,c)=(\partial a+hb,-\partial b+kc,\partial c),
\]
where $hk=0$.  The only possible higher differential is
\[
 d^2:E^2_{4,q}\longrightarrow E^2_{2,q+1}.
\]
If $\partial c=0$ and $\partial b=kc$, it is represented by $hb$.
Thus a differential lowering core size by at most one can still produce
a nonzero secondary differential.

An explicit eight-state example makes this bound sharp.  Take
$Q=A_2\oplus R\oplus A_2$, with bottom antichain $\{a,a'\}$,
top antichain $\{y,y'\}$, and
$R=\{v,b,w,x\}$ with relations $v<x$, $b<w$, $b<x$.
Make $a,b$ one-only and $x,y$ zero-only; split the other states.
The only nonzero first-page groups are
$E^1_{2,1}=\mathbb Z$ and $E^1_{4,0}=\mathbb Z$.
The four-core link consists of $b_1,x_0$; all triple links are
contractible; only the edge $ay$ has a nontrivial available link, a circle
up to homotopy.  The paths
\[
 x_0-v_0-v_1-a_1-b_1,\qquad
 x_0-y_0-w_0-w_1-b_1
\]
in the two relevant triple links fill the four-core difference.
Their difference generates the sole loop in the $ay$ link.  Therefore
$d^2$ is a unit $\mathbb Z\to\mathbb Z$, and $E^3=0$.
The supplementary signed-chain certificate verifies these fillings and
the integral homology directly.

Nor do Ferrers or chordal-bipartite graphs force $E^2$ collapse.
The seven-state poset with covers
$0<1$, $0<2$, $1<3$, $2<3$, $2<5$, $3<4$, $3<6$, $5<6$,
and profile $(1,2,0,0,2,2,0)$ has exclusive graph $K_{1,3}$.
Its integral $E^1=E^2$ has $\mathbb Z^2$ in bidegree $(2,1)$
and $\mathbb Z$ in $(4,0)$; $d^2$ is a primitive injection.
Its only nonzero cone group is $H_2=\mathbb Z$.
These are explicit counterexamples, not global minimality claims.

\subsection{An unbounded interacting class and new torsion}
\begin{proposition}[All-exclusive star]
Let $R$ be a nonempty finite poset.  Take $Q=\{p\}\oplus R$, make
$p$ one-only and every state of $R$ zero-only.  Then $E^2=E^\infty$
integrally and
\[
 \widetilde H_n(\operatorname{Cone}(\Delta\pi);\mathbb Z)
 \cong H_{n-1}(\Delta R;\mathbb Z),
\]
where the right side is ordinary homology.
\end{proposition}
\begin{proof}
Every core is $\{p_0\}\cup\sigma_1$ for a nonempty chain $\sigma$
in $R$.  Its available link is the isolated vertex $p_1$ disjoint from
the cone $\Delta(R_{\le\min\sigma}\times\{0\})$.
The difference-of-components generators identify all nonzero link
homology with $\mathbb Z$ in degree zero and identify their inclusion
maps with the identity.  Consequently $d^1$ is, up to a uniform sign,
the ordinary simplicial boundary of $\Delta R$.  The single row
collapses on $E^2$ and has only one filtration grade in each total degree.
\end{proof}

For the nonempty-face poset of a triangulated projective plane, all these
available links have torsion-free homology, and every coarse edge link
is a vertex link in its barycentric subdivision, hence a circle.
Nevertheless the interaction differential creates $\mathbb Z/2$ in
cone degree two.  The explicit six-vertex triangulation in the
supplement gives a 32-state example, with every link and the reduced
interaction complex checked by integral Smith normal form.

\begin{remark}[Scope and priority]
Filtered-complex spectral sequences, link/cofiber recursions, and their
extension issues are classical.  The complexes $\Lambda_S$ are precisely
vertex-cover obstruction complexes in the sense of
Chach\'olski--Jin--Scolamiero--Tombari~\cite{ChacholskiEtAl2021}.
The results here explicitly identify which such complexes interact in
the observation product model, with positive class restrictions and
integral obstructions.  They do not introduce a new general obstruction
theory or supersede poset-fiber theorems.
\end{remark}
